% ============================================================================
%  Chapter 16 -- Permanent Memory                                (Lecture 16)
% ============================================================================
\chapter{Permanent Memory}
\label{ch:storage}

\epigraphline{\ldots a datum does not exist unless it has been saved\ldots}
  {Prof.\ Almerico Murli, 1990s}

\begin{chapgoals}
By the end of this chapter you will be able to
\begin{itemize}
\item name the functional units of a computer and recognise them on a motherboard and on the \prg kit;
\item distinguish ROM, PROM, EPROM, EEPROM and flash, and say what firmware is;
\item explain why a cache makes a computer faster, and what a cache miss is;
\item describe tracks, sectors and blocks of a magnetic disk and the components of its access time;
\item explain how an optical disk stores bits, and compare CD and DVD;
\item place flash, solid-state disks, hard disks and tapes in a multi-level storage system.
\end{itemize}
\end{chapgoals}

% ----------------------------------------------------------------------------
\section{Warm-up on the machine}

\asmfile[caption={Switch the kit off and on again.},label=lst:st:warm]{code/l16/c03.S}

A game keeps the best score in a variable (Listing~\ref{lst:st:warm}). After
the kit has been switched off and on again the title is still there, but the
best score is 0. The reason is in the two section directives. The string is
in \code{.rodata}, which is part of the cartridge file, and the cartridge is a
file in \emph{flash} memory, which is non-volatile. The variable is in
\code{.data}, which is copied into RAM at the start and lives in SRAM, which
is volatile.

To keep a score for ever, it must be written to a permanent memory. Which
ones exist? (On \prg, persistent scores are a service of the Cartridge Store;
the point here is where each section of a program lives.)

% ----------------------------------------------------------------------------
\section{The components of a computer}

\subsection{Functional units}

\begin{description}[style=nextline,leftmargin=1.5em]
\item[Processor.] Executes the instructions. On the kit: the \riscv core of
  the ESP32-C6.
\item[Main memory.] Holds the running programs and their data. On the kit:
  512~KB of SRAM inside the chip.
\item[Secondary memory.] Keeps programs and data when the power is off. On
  the kit: the flash chip that holds the firmware and the cartridges.
\item[Input/output devices.] Video card and monitor, keyboard, mouse, printer,
  storage. On the kit: display, joystick, amplifier.
\item[Communication lines.] The buses: data and address lines shared by all
  the units.
\end{description}

A personal computer spreads these units over a motherboard; a microcontroller
puts almost all of them in one package.

\subsection{The motherboard}

A \term{motherboard} (Figure~\ref{fig:st:board}) carries the sockets for one
or more processors and, close to them, the slots for the memory modules; a
\emph{chipset} that connects processor, memory and peripherals; expansion
buses, today PCI Express slots for video cards and other boards; connections
for storage (Serial ATA, M.2) and for the outside world (USB, Ethernet,
audio); and a flash chip with the firmware that starts the machine. From an
architectural point of view a board of 2002 and one of today have much in
common.

\begin{figure}[htb]
\centering
\includegraphics[width=0.7\textwidth]{workstation_board.jpg}
\caption{A dual-processor workstation motherboard.}
\label{fig:st:board}
\end{figure}

\subsection{Firmware: the BIOS and its heirs}

\begin{definition}{Firmware}
The \term{firmware} is the basic software of a machine, kept in a non-volatile
memory and used by the processor from the moment in which it is switched on.
\end{definition}

In a personal computer it was called the BIOS: a program in a non-volatile
chip, specific to a given motherboard. Its instructions and data are divided
into four parts: the \emph{POST} (power-on self test), which checks that the
devices work; the \emph{SETUP}, the configuration of the devices; the
\emph{BOOT}, which starts the operating system; and the \emph{DRIVERS}, basic
low-level functions for the keyboard, the disks and other devices. Today's
PCs use UEFI firmware, which does the same jobs.

\begin{prgmachine}
On the ESP32-C6 a boot program in ROM starts a bootloader in flash, which
starts the \prg firmware, which loads the cartridge chosen from its
\emph{Setup} menu. Each stage finds, checks
and starts the next: a chain that goes from the first instruction after reset,
through layers of permanent memory, to the first instruction of the reader's
\code{init}.
\end{prgmachine}

% ----------------------------------------------------------------------------
\section{Non-volatile memories}

\subsection{From ROM to flash}

Semiconductor memories that keep their content without power have evolved
towards an ever easier way of writing them (Table~\ref{tab:st:rom}). Reading
them takes a time of the order of nanoseconds.

\begin{table}[htb]
\centering\small
\caption{From ROM to flash.}\label{tab:st:rom}
\begin{tabularx}{\textwidth}{@{}lP{0.24\textwidth}LL@{}}
\toprule
\textbf{Memory} & \textbf{Written by} & \textbf{Erased} & \textbf{Typical use}\\
\midrule
ROM & the manufacturer, once and for all & never & start-up code, old game cartridges\\
PROM & the user, once & never & small production runs\\
EPROM & the user, many times & all at once, with ultraviolet light through a quartz window (about 15 minutes) & development, BIOS of early PCs\\
EEPROM & the user, in place & electrically, one byte at a time & configuration data\\
Flash & the user, in place & electrically, in blocks & firmware, USB keys, memory cards, solid-state disks\\
\bottomrule
\end{tabularx}
\end{table}

The acronyms are \emph{read-only memory}, \emph{programmable} ROM,
\emph{erasable} PROM and \emph{electrically erasable} PROM. A ROM is, from
the logical point of view, a combinational machine: a table that associates a
fixed word with each address, as the tables \code{dx} and \code{dy} of
\vermone do.

\subsection{Flash memory}

Flash is the most recent kind of EEPROM: it can be reprogrammed in place,
without being removed from the board, and it is today the cheapest
semiconductor memory. It has three peculiarities.

\begin{itemize}
\item It is erased by \emph{blocks} of bytes, not one byte at a time. Erasing
  sets all the bits of a block to 1, and writing can only turn 1s into 0s.
\item Each block survives a limited number of erase cycles, so that the
  controller spreads the writes evenly over the blocks.
\item Reading is fast; writing and, above all, erasing are slow.
\end{itemize}

The flash of the kit holds the firmware and the cartridges: uploading a
cartridge erases and rewrites some of its blocks. And this is why the best
score of the warm-up is not simply written back into the cartridge at every
frame: flash is not a place for variables.

\begin{figure}[htb]
\centering
\begin{subfigure}[b]{0.3\textwidth}\centering\includegraphics[width=\textwidth,height=34mm,keepaspectratio]{eprom.jpg}\caption{An EPROM with its quartz window}\end{subfigure}\hfill
\begin{subfigure}[b]{0.3\textwidth}\centering\includegraphics[width=\textwidth,height=34mm,keepaspectratio]{hard_disk.jpg}\caption{A 6 TB hard disk drive}\label{fig:st:hdd}\end{subfigure}\hfill
\begin{subfigure}[b]{0.3\textwidth}\centering\includegraphics[width=\textwidth,height=34mm,keepaspectratio]{ssd.jpg}\caption{A solid-state disk}\end{subfigure}
\caption{Three permanent memories.}
\end{figure}

% ----------------------------------------------------------------------------
\section{The cache}

\subsection{Why a cache?}

During the clock cycles needed to read or to write the memory the processor
waits; the slower the memory, the more cycles are lost. Why not build
memories as fast as the registers of the processor? Because they would have
to be on the same chip, since a connection through a bus is slower; and then
they could not be very large, or the chip would grow in size and in cost. The
compromise is a small, fast memory on the chip of the processor, and one or
more larger and slower memories reached through the bus.

\subsection{How a cache works}

\begin{definition}{Cache}
A \term{cache} (from the French \emph{cacher}, to hide) is a memory of small
capacity but fast, placed between the processor and the main memory, which
keeps a copy of part of the content of the latter.
\end{definition}

The processor looks in the cache first (Figure~\ref{fig:st:cache}). There is
a \term{hit} if the content of the requested address is there, and it arrives
at once. There is a \term{miss} if it is not, and then it must be copied from
the RAM into the cache; a fixed number of neighbouring locations is copied
with it. The program does not stop and does not notice: the access is simply
slower.

\begin{figure}[htb]
\centering
\begin{tikzpicture}
  \node[blkacc,minimum width=36mm,minimum height=9mm] (p) at (0,0) {processor};
  \node[blkamb,minimum width=36mm,minimum height=11mm] (c) at (0,-1.8) {\textbf{cache}\\small, fast};
  \node[blkvio,minimum width=36mm,minimum height=13mm] (m) at (0,-3.9) {\textbf{main memory}\\large, slow};
  \draw[<->,bus] (p) -- (c) node[midway,right,lbl]{every access};
  \draw[<->,bus] (c) -- (m) node[midway,right,lbl]{only on a miss: a block};
\end{tikzpicture}
\caption{A cache between the processor and the main memory.}
\label{fig:st:cache}
\end{figure}

Why copy the neighbours? Because of an empirical property of programs.

\begin{prgconcept}[The principle of locality]
A program that has just used an address will probably use it again soon
(\emph{temporal locality}) and will probably use its neighbours
(\emph{spatial locality}). Loops and arrays are the reason.
\end{prgconcept}

The more hits, the shorter the average access time. If a fraction $h$ of the
accesses are hits, the cache answers in a time $t_c$ and a miss costs in
addition the time $t_m$ of the memory, the average is
\[ t = h\,t_c+(1-h)\,(t_c+t_m). \]
With $t_c=1$~ns, $t_m=50$~ns and 95\% of hits,
$t=0.95\times 1+0.05\times 51=3.5$~ns: the system behaves almost as if the
whole memory were as fast as the cache.

\subsection{Cache levels and typical machines}

Real processors have several levels of cache, all built with SRAM
(Table~\ref{tab:st:levels}). The first level, internal, keeps a continuous
flow inside the processor; the second limits the slowdown due to the main
memory; the third, of 8 to 32~MB and more, is shared by the cores of the same
processor.

\begin{table}[htb]
\centering\small
\caption{Cache levels on three machines.}\label{tab:st:levels}
\begin{tabularx}{\textwidth}{@{}lllL@{}}
\toprule
 & \textbf{PC of 2002} & \textbf{PC of 2020} & \textbf{ESP32-C6 (the kit)}\\
\midrule
Processor & 1 core & 4 cores & 1 \riscv core, 160 MHz\\
Cache L1 & 64 KB & 64 KB per core & none for the SRAM\\
Cache L2 & 512 KB & 512 KB per core & none\\
Cache L3 & none & 8 MB, shared by the cores & none\\
Main memory & 256 MB & 16 GB & 512 KB of SRAM\\
Storage & 40 GB hard disk & 1 TB + 1 TB, mirrored & a few MB of flash, read through a small cache\\
\bottomrule
\end{tabularx}
\end{table}

On the ESP32-C6 the SRAM is already as fast as the core needs; only the
external flash sits behind a cache.

% ----------------------------------------------------------------------------
\section{Secondary memory}

\subsection{Characteristics}

The \term{secondary memory}, or \term{mass storage}, keeps all the programs
and data of a computer, whereas the main memory holds the programs being
executed and the data that they need. Compared with the main memory it is

\begin{itemize}
\item non-volatile: the data are not lost when the computer is switched off;
\item of large capacity: several orders of magnitude more;
\item of low cost per bit: several orders of magnitude less;
\item slower: by some orders of magnitude.
\end{itemize}

Its information is organised in \term{blocks}, usually of a kilobyte or more,
not in single bytes. Some media allow \emph{direct} access to any block;
others only \emph{sequential} access.

\subsection{Magnetic disks: tracks, sectors, blocks}

A magnetic disk (Figure~\ref{fig:st:hdd}) is made of platters coated with a
magnetisable material and kept spinning at constant speed, and of heads that
read and write the information. The surface is divided into concentric
\term{tracks} and into \term{sectors}, slices of the disk
(Figure~\ref{fig:st:disk}). A sector cuts each track into blocks: the block
is the smallest addressable unit, identified by the number of its track and
of its sector. \term{Formatting} a disk means dividing its surface into
tracks and sectors. Access is direct: the head can be positioned on any block
without passing over all the previous ones. Several sectors form a
\emph{cluster}; the tracks at the same radius on all the platters form a
\emph{cylinder}.

\begin{figure}[htb]
\centering
\begin{tikzpicture}
  \fill[prgpanel] (0,0) circle (2.6);
  \fill[prgamber!45] (0,0) -- (30:2.6) arc (30:60:2.6) -- cycle;
  \begin{scope}
    \clip (0,0) -- (30:2.7) arc (30:60:2.7) -- cycle;
    \fill[prgred!70,even odd rule] (0,0) circle (1.9) (0,0) circle (1.5);
  \end{scope}
  \draw[prgaccent,line width=2pt] (0,0) circle (1.7);
  \foreach \r in {0.7,1.1,1.5,1.9,2.3,2.6} \draw[prgsteel,thin] (0,0) circle (\r);
  \foreach \a in {0,30,...,330} \draw[prgsteel,thin] (\a:0.7) -- (\a:2.6);
  \fill[white] (0,0) circle (0.7); \draw[prgsteel] (0,0) circle (0.7); \fill[prgsteel] (0,0) circle (0.12);
  \draw[<-,thick] (150:1.7) -- (-3.6,1.6) node[left,lbl,text=prgaccent]{a track};
  \draw[<-,thick] (45:2.4) -- (3.4,2.2) node[right,lbl,text=prgamber]{a sector};
  \draw[<-,thick] (45:1.7) -- (3.4,0.6) node[right,lbl,text=prgred]{a block};
  \draw[line width=3pt,prgink] (3.3,-2.4) -- (-10:1.75);
  \fill[prgink] (-10:1.75) circle (0.1);
  \node[lbl,anchor=west] at (3.4,-2.4) {arm and head};
\end{tikzpicture}
\caption{The surface of a magnetic disk.}
\label{fig:st:disk}
\end{figure}

\subsection{Access time of a magnetic disk}

To read or to write a block the head must be brought to the right address;
during the transfer the head is still and the disk moves. The access time is
the sum of three components.

\begin{description}[style=nextline,leftmargin=1.5em]
\item[Seek time.] The time to move the head radially to the right cylinder,
  and so to the right track.
\item[Latency.] The time to wait for the desired sector to pass under the
  head; on average, half a revolution.
\item[Transfer time.] The time for the block to pass under the head.
\end{description}

A disk that turns at 7\,200 revolutions per minute takes $60/7200 = 8.33$~ms
per revolution, and its average latency is about 4.2~ms. Hard disks are
sealed against dust, have capacities of terabytes and access times of
milliseconds. A millisecond is a million nanoseconds: during one access to a
disk a processor at 160~MHz could execute more than a hundred thousand
instructions.

\subsection{Hard disks and solid-state disks}

A \term{hard disk drive} (HDD) has magnetic platters and moving heads; its
access time, of milliseconds, is dominated by mechanics; it has the lowest
cost per terabyte and is still the choice for large archives and for backups.
A \term{solid-state disk} (SSD) is flash memory with a controller, and has no
moving parts; its access time is of tens of microseconds; it is connected
through SATA or, faster, through NVMe on PCI Express. It is the standard
storage of laptops and telephones, and is the same technology as the flash
chip of the kit.

% ----------------------------------------------------------------------------
\section{Optical disks}

\subsection{Bits read with light}

An optical disk is a rigid disk on which a laser beam of a few milliwatts is
projected; the reflection changes between the \term{pits} and the surface,
the \term{land} (Figure~\ref{fig:st:pits}).

\begin{prgwarn}
0 and 1 are \emph{not} pit and land. A 1 is a \emph{change}, from pit to land
or from land to pit; a 0 is no change.
\end{prgwarn}

Encoding a 1 as a transition is an idea used in many serial links as well,
because it lets the receiver recover the rhythm of the bits from the signal
itself. Optical disks are usually written once, because the physical changes
are irreversible: they are used to distribute programs and as archives that
must not be modified. In a CD-Recordable, two layers of polycarbonate enclose
a thin reflective sheet covered by a translucent dye, and recording burns
points in the dye; rewritable disks use a material that can change back.

\begin{figure}[htb]
\centering
\begin{subfigure}[b]{0.3\textwidth}\centering\includegraphics[width=\textwidth,height=34mm,keepaspectratio]{pit_land.jpg}\caption{Pits and lands}\label{fig:st:pits}\end{subfigure}\hfill
\begin{subfigure}[b]{0.3\textwidth}\centering\includegraphics[width=\textwidth,height=34mm,keepaspectratio]{cdrom_drive.jpg}\caption{Inside a CD-ROM drive}\end{subfigure}\hfill
\begin{subfigure}[b]{0.3\textwidth}\centering\includegraphics[width=\textwidth,height=34mm,keepaspectratio]{dvd_drive.jpg}\caption{Inside a DVD drive}\end{subfigure}
\caption{Optical storage.}
\end{figure}

\subsection{The compact disc}

On a compact disc the information is stored on a single spiral track, with
constant recording density; the distance between two successive turns of the
spiral, the \emph{track pitch}, is 1.6 micrometres. Two ways of turning the
disk are possible. With \emph{constant linear velocity} (CLV) the disk turns
faster when the head reads near the centre, from 200 to 500 revolutions per
minute, so that the track passes under the head always at the same speed
(1.2~m/s at the basic speed, set by the bit rate needed to play music). With
\emph{constant angular velocity} (CAV) the revolutions are constant and so is
the number of bits per sector, with the result that the outer surface is
under-used.

\subsection{DVD and CD compared}

The DVD (\emph{digital versatile disk}) has smaller pits and a tighter spiral,
and up to two layers and two sides (Table~\ref{tab:st:dvd}).

\begin{table}[htb]
\centering\small
\caption{DVD and CD.}\label{tab:st:dvd}
\begin{tabular}{@{}lll@{}}
\toprule
 & \textbf{DVD} & \textbf{CD}\\
\midrule
Disk diameter & 120 mm & 120 mm\\
Disk thickness & 1.2 mm ($0.6$ mm $\times$ 2) & 1.2 mm\\
Surfaces / layers & 1 or 2 / 1 or 2 & 1 / 1\\
Minimum pit size & 0.4 $\mu$m & 0.83 $\mu$m\\
Track pitch & 0.74 $\mu$m & 1.6 $\mu$m\\
Average transfer rate & 4.7 MB/s & 0.15 MB/s\\
Capacity, 1 layer, 1 side & about 5 GB & about 0.68 GB\\
Capacity, 2 layers, 2 sides & 17 GB & --\\
\bottomrule
\end{tabular}
\end{table}

% ----------------------------------------------------------------------------
\section{Storage on several levels}

A modern computing system has several levels of permanent memory
(Figure~\ref{fig:st:levels}), often managed automatically: the data that are
used often migrate towards the fast levels, and the rest sinks towards the
cheap ones. It is the idea of the cache, applied to storage; and from a chip
of a few megabytes on the kit to rooms of petabytes
(Figure~\ref{fig:st:racks}) the hierarchy is the same at every scale.

\begin{figure}[htb]
\centering
\begin{tikzpicture}[node distance=4mm]
  \node[blkred,text width=21mm,minimum height=17mm] (a) {\textbf{RAM disk}\\[1pt]\footnotesize fastest, volatile};
  \node[blkamb,text width=21mm,minimum height=17mm,right=of a] (b) {\textbf{Solid-state disk}\\[1pt]\footnotesize working storage};
  \node[blkacc,text width=21mm,minimum height=17mm,right=of b] (c) {\textbf{Hard disk}\\[1pt]\footnotesize large, cheap};
  \node[blk,text width=21mm,minimum height=17mm,right=of c] (d) {\textbf{Disk library}\\[1pt]\footnotesize on when needed};
  \node[blkvio,text width=21mm,minimum height=17mm,right=of d] (e) {\textbf{Tape library}\\[1pt]\footnotesize archives};
  \foreach \x/\y in {a/b,b/c,c/d,d/e} \draw[<->,thick] (\x) -- (\y);
\end{tikzpicture}
\caption{A multi-level permanent memory.}
\label{fig:st:levels}
\end{figure}

\begin{figure}[htb]
\centering
\includegraphics[width=0.6\textwidth]{storage_racks.jpg}
\caption{Disk and tape libraries in a data centre.}
\label{fig:st:racks}
\end{figure}

% ----------------------------------------------------------------------------
% Generated by tools/textbook/export_listings.js from decks/plan.js and code/l16.S: do not edit.
\section{The program of this chapter}\label{sec:lab:l16}

The listings of this chapter are parts of one program, \fname{code/l16.S} in the
repository of the course, printed in full in Listing~\ref{lst:lab:l16}. The
automatic tests of the course material build it and run it on the headless
runner of \prgqt{}, with scripted input, and check the frame;
Figure~\ref{fig:lab:l16} is the frame that the test wrote.

\begin{description}[style=nextline,leftmargin=1.5em]
\item[On the screen.] The title HIGH SCORE and two bars: the score of this run, which grows while A is held and falls back to zero when A is released, and the best score, which never goes down. Restart the cartridge: the title is still there, the best score is zero again.
\item[What to press.] Hold A, release it, then restart the cartridge.
\end{description}

\begin{figure}[htb]
\centering
\includegraphics[width=0.5\textwidth]{lab_l16.png}
\caption{\fname{l16.S} on \prgqt{}.}\label{fig:lab:l16}
\end{figure}

\asmfile[caption={The complete program of the chapter, \fname{code/l16.S}.},label=lst:lab:l16]{code/l16/program.S}

To run it on \prgqt{} the commands are those of Listing~\ref{lst:lab:l16:run};
Section~\ref{sec:prg:testing} gives the same steps for QEMU and for the board.

\begin{shell}[caption={Building \fname{l16.S} and running it on \prgqt{}.},label=lst:lab:l16:run]
# 1. start PRG32-QT on this computer (it listens on port 8080)

# 2. in the PRG32 folder, with ESP-IDF active: build the cartridge
python3 -m prg32 cartridge build ~/arc12/code/l16.S --portable \
    --entry-prefix l16 --name l16 --out l16.prg32

# 3. send it to PRG32-QT and start it
python3 -m prg32 esp32c6 upload-and-run l16.prg32 \
    --url http://localhost:8080 --slot cart0

# 4. look, press, compare with the previous slide
\end{shell}


% ----------------------------------------------------------------------------
\section{Check your understanding}

\begin{quiz}{Which memory can be erased electrically, in blocks, without removing it from the board?}
\item ROM
\item PROM
\item EPROM
\item Flash
\end{quiz}

\begin{quiz}{What happens on a cache miss?}
\item The program stops with an error
\item The datum is fetched from main memory and copied into the cache, with its neighbours
\item The cache is erased
\item The processor skips the instruction
\end{quiz}

\begin{quiz}{Which component of the access time of a hard disk depends on the rotation speed?}
\item Seek time
\item Latency
\item Formatting time
\item Boot time
\end{quiz}

\begin{quiz}{On an optical disk, a 1 is encoded as:}
\item a pit
\item a land
\item a transition between pit and land
\item a double pit
\end{quiz}

\section{Exercises}

\begin{exercise}[paper and \prg]
\begin{enumerate}[nosep]
\item A disk turns at 5\,400 revolutions per minute: what is its average
  latency?
\item A cache answers in 2~ns and the memory in 80~ns: what is the average
  access time with 90\% of hits? And with 99\%?
\item Say where each of these lives on the \prg kit: the boot code, the
  firmware, a cartridge, the variable \code{counter}, the register \reg{t0}.
\item Find the specifications of the storage of your laptop or telephone:
  technology, capacity, interface.
\item On \prg, time two loops that sum a large array: one in order, one that
  jumps by a large step. Is there a difference, and why might there not be on
  this chip?
\end{enumerate}
\expected{Latency: half of $60/5400=11.1$~ms, that is, about 5.6~ms. Access
time: $0.9\times 2+0.1\times 82=10$~ns, and $0.99\times 2+0.01\times 82=2.8$~ns.
The five things live in ROM, in flash, in flash (or in RAM once loaded), in
SRAM, and in a register of flip-flops inside the processor. In the last task
the SRAM of the ESP32-C6 has no cache in front of it, so that the pattern of
access matters little; on a laptop the difference is large.}
\end{exercise}

\section{Glossary}

\begin{glossarytable}
motherboard & scheda madre\\
firmware & firmware\\
non-volatile & non volatile\\
cache hit / miss & successo / fallimento di cache\\
mass storage & memoria di massa\\
track / sector / block & traccia / settore / blocco\\
seek time & tempo di posizionamento\\
latency & tempo di latenza\\
solid-state disk & disco a stato solido\\
tape library & libreria a nastri\\
\end{glossarytable}

\section{Conclusions and reading}

Non-volatile memories usually hold the basic software of a machine, the
firmware. The cache improves performance by keeping a copy of part of the
main memory close to the processor, and the same idea, a small fast level in
front of a large slow one, is used again in storage systems. Permanent
memories can be semiconductor (flash), magnetic or optical, and a modern
system has several levels of them, often managed automatically.

\paragraph{Reading.} \citet{patterson2020}, sections 5.1--5.4;
\citet{stallings2021}, the chapters on cache and on external memory;
\citet{tanenbaum2013}, chapter~2.

The third block is complete. The next chapter uses it in the third part of
\vermone: a body, rules and states.
